\clearpage
\VerdierPhysicalPageAnchor{0226}%
\begingroup\small\setlength{\parskip}{0pt}\setlength{\emergencystretch}{1.5em}
\setlength{\abovedisplayskip}{4pt plus 2pt minus 2pt}%
\setlength{\belowdisplayskip}{4pt plus 2pt minus 2pt}%

\VerdierSemanticAnchor{ch4d-4-1-2-cont}
qui définissent une filtration du type $\widetilde{\mathbb Z}$ de $Y$
qu'on appelle \emph{première filtration canonique de $Y$}, et par suite un
foncteur, que nous notons $Y$, de $\widetilde{\mathbb Z}$ dans
$\operatorname{comp}(\mathcal A)$. On en déduit par la construction
rappelée en (4.1.1) un objet spectral à valeurs dans $K(\mathcal A)$ :
$(\mathcal Y(p,q),\delta(p,q,r))$. On se propose de décrire un objet
spectral isomorphe à $(\mathcal Y(p,q),\delta(p,q,r))$. Posons tout d'abord :

\VerdierSemanticAnchor{ch4d-4-1-2-3}
\begin{equation}
\resizebox{0.98\textwidth}{!}{%
$\displaystyle
\begin{aligned}
\text{a) }& -\infty<p<q<+\infty:\\[-2pt]
&Y(p,q)=\cdots\longrightarrow0\longrightarrow0\longrightarrow Y^{-q}
\longrightarrow Y^{-q+1}\longrightarrow\cdots\longrightarrow Y^{-p-2}
\longrightarrow Y^{-p-1}\longrightarrow0\longrightarrow0
\longrightarrow\cdots;\\[2pt]
\text{b) }& -\infty<p<+\infty:\\[-2pt]
&Y(p,+\infty)=\cdots\longrightarrow Y^i\longrightarrow Y^{i+1}
\longrightarrow\cdots\longrightarrow Y^{-p-2}\longrightarrow Y^{-p-1}
\longrightarrow0\longrightarrow0\longrightarrow\cdots;\\[2pt]
\text{c) }&Y(-\infty,q)=Y(q)\quad(q\in\widetilde{\mathbb Z}),\qquad
Y(p,p)=0\quad(p\in\widetilde{\mathbb Z}).
\end{aligned}$% 
}
\tag{4.1.2.3}
\end{equation}
Les différentielles de ces complexes sont induites par les différentielles
du complexe $Y$.

Désignons par :

\VerdierSemanticAnchor{ch4d-4-1-2-4}
\begin{equation}
\varphi\bigl((p,q),(p',q')\bigr)\colon
Y(p,q)\longrightarrow Y(p',q'),\qquad p\leq p',\quad q\leq q',
\tag{4.1.2.4}
\end{equation}
la classe d'homotopie du morphisme de complexes :

\VerdierSemanticAnchor{ch4d-4-1-2-5}
\begin{equation}
\psi\bigl((p,q),(p',q')\bigr)\colon
Y(p,q)\longrightarrow Y(p',q'),
\tag{4.1.2.5}
\end{equation}
dont les composantes sont des morphismes identiques quand la source et le
but sont des objets non nuls. On vérifie immédiatement qu'on a ainsi défini
un foncteur de $\mathcal{F}\ell_1(\widetilde{\mathbb Z})$ dans
$K(\mathcal A)$.

Soit $(p,q,r)$ un élément de
$\mathcal{F}\ell_2(\widetilde{\mathbb Z})$. Désignons par :

\VerdierSemanticAnchor{ch4d-4-1-2-6}
\begin{equation}
\delta(p,q,r)\colon Y(q,r)\longrightarrow Y(p,q)[1]
\tag{4.1.2.6}
\end{equation}
la classe d'homotopie du morphisme de complexes :

\VerdierSemanticAnchor{ch4d-4-1-2-7}
\begin{equation}
d(p,q,r)\colon Y(q,r)\longrightarrow Y(p,q)[1],
\tag{4.1.2.7}
\end{equation}
dont les composantes sont :

\VerdierSemanticAnchor{ch4d-4-1-2-8}
\begin{equation}
\begin{aligned}
d(p,q,r)^i&=0,\qquad i\neq-q-1,\\
d(p,q,r)^{-q-1}&=-d_Y^{-q-1}.
\end{aligned}
\tag{4.1.2.8}
\end{equation}
\endgroup
